\section{Hardware Abstraction and Data Access}

\begin{quotation}
    \textit{\textbf{Author's note:} This section will be MORE technical.\\
    For students who are studyng with my notes, I suggest to read it carefully.\\
    Moreover, I \underline{think} (PERSONAL OPINION) that it is just an overview and it is not necessary to momorize all the details, eg. the command line examples.} 
\end{quotation}

\noindent
In the context of a "sound and admissible" investigation, Prof. Atzeni highlights the critical transition from high-level OS to low-level hardware interaction.

\subsection{The File System and Mounting Risks}
OS provides a simplified interface (files and folders)to interact with storage, this abstraction presents specific forensic risks during the \textbf{mounting} process:
\begin{itemize}
    \item \textbf{Automatic Alteration:} modern OS may automatically update metadata (e.g., last access time) or attempt to repair file system inconsistencies upon mounting, potentially destroying evidence.
    \item \textbf{Forensic Precaution:} an investigator must ensure the device is mounted in \textbf{read-only} mode at both the file system and block device levels to maintain data integrity.
\end{itemize}

\subsection{Unix Device Taxonomy}
Under Unix-derived systems, hardware is abstracted as files within the \texttt{/dev} directory.

\begin{table}[h!]
\centering
\small
\begin{tabular}{|l|p{6cm}|l|}
\hline
\textbf{Device Type} & \textbf{Characteristics} & \textbf{Examples} \\ \hline
\textbf{Block Device} & Handles data in fixed-size blocks; allows \textbf{Random Access}. & HDD, SSD, USB \\ \hline
\textbf{Character Device} & Handles data as a \textbf{continuous stream}; no random access. & Terminals, Consoles \\ \hline
\textbf{Network Device} & Managed via packets and sockets using system calls. & \texttt{eth0}, \texttt{wlan0} \\ \hline
\textbf{Pseudo-device} & Software interfaces emulating hardware features. & \texttt{/dev/null}, \texttt{/dev/loop} \\ \hline
\textbf{Special-purpose} & Used for kernel interaction and memory management. & \texttt{/dev/mem}, \texttt{/dev/kmem} \\ \hline
\end{tabular}
\caption{Classification of devices according to Prof. Atzeni.}
\end{table}

\noindent
The distinction between how data is accessed is fundamental for forensic "bit-by-bit" acquisition.

\subsubsection*{1. Access via File System (Logical)}
Access occurs through the \textbf{mount point} using standard OS tools (e.g., \texttt{ls}, \texttt{cp}).
\begin{itemize}
    \item \textbf{Pro:} Easy to navigate; metadata is managed by the OS.
    \item \textbf{Cons:} Limited to active files; risk of metadata corruption; cannot see "unallocated" space or deleted files.
\end{itemize}

\subsubsection*{2. Direct Access (Physical/Raw)}
Access occurs directly at the \textbf{block device} level (e.g., \texttt{/dev/sda1}).
\begin{itemize}
    \item \textbf{Requirement:} Requires \textbf{root/administrator} privileges.
    \item \textbf{Functionality:} Treats the storage as unstructured blocks of data.
    \item \textbf{Forensic Importance:} Essential for creating bit-by-bit copies and performing "data carving" (recovering deleted information) using tools like \texttt{dd} or \texttt{hexdump}.
\end{itemize}



\subsubsection{Unix Forensic Commands}
\begin{itemize}
    \item \texttt{lsblk / fdisk -l}: Identification of connected devices regardless of mounting status.
    \item \texttt{mount -o ro,noexec,nosuid}: Secure mounting parameters.
    \begin{itemize}
        \item \textbf{ro}: Read-only access.
        \item \textbf{noexec}: Prevents execution of binaries.
        \item \textbf{nosuid}: Disables set-user-ID bits (prevents privilege elevation).
    \end{itemize}
    \item \texttt{dmesg}: Reviewing kernel messages to identify newly attached hardware.
\end{itemize}


\subsection{Procedures regading File System }

A forensically sound environment requires a multi-layered write-protection strategy to prevent any data alteration:
\begin{enumerate}
    \item \textbf{Hardware Level:} Use of physical hardware write blockers to intercept write signals.
    \item \textbf{Device Level (IOCTL):} Setting the block device to read-only using the \texttt{blockdev} utility.
    \begin{itemize}
        \item \texttt{blockdev --setro /dev/sdb1} (Set Read-Only)
        \item \texttt{blockdev --getro /dev/sdb1} (Verify status)
    \end{itemize}
    \item \textbf{File System Level:} Mounting the partition with the \texttt{-o ro} flag.
\end{enumerate}
\textit{ Note:} Device-level protection overrides file system requests, ensuring that even accidental write attempts by the OS fail.

\subsubsection{Disk Images}
A \textbf{disk image} is a file containing a complete, bit-per-bit replica of a storage medium (e.g., created via \texttt{dd}).
\begin{itemize}
    \item Images contain everything: file structures, metadata (MFT, journals), and unallocated space.
    \item Tools like \texttt{imount} or CAINE's \textit{UnBlock} allow investigators to mount these images and interact with them as if they were physical drives.
\end{itemize}

\subsubsection{File Attributes}
The OS uses several attributes to manage files, but an investigator must look beyond the surface:
\begin{itemize}
    \item \textbf{Magic Numbers:} The first few bytes of a file that identify its true type, regardless of the filename extension. This is critical to detect renamed malicious files.
    \item \textbf{MAC Times:} Metadata tracking Modification, Access, and Change (creation) times, essential for timeline reconstruction.
    \item \textbf{Size vs. Physical Space:} Files are stored in blocks; accessing the device directly (\textit{Direct Access}) allows an investigator to see data hidden in "slack space" or within sub-block fragments.
\end{itemize}